% =====================================================================
%  CORPO del volume "Soluzioni" (comune alle versioni light e dark).
%  Quando aggiungi un modulo: decommenta la sua riga \input{solNN}.
%  I capitoli vanno inclusi TUTTI e in ordine, così la numerazione
%  (Capitolo 3 = Modulo 3, "Codice 3.5", ...) coincide con la dispensa.
% =====================================================================
\begin{titlepage}
  \centering
  \vspace*{3cm}
  {\Huge\bfseries Programmazione 2\par}
  \vspace{0.6cm}
  {\Large Soluzioni degli esercizi al PC\par}
  \vspace{1.5cm}
  \begin{tcolorbox}[colback=blue!4!sfondo,colframe=blue!55!black,width=0.85\textwidth,title=Come usare questo volume]
    Qui trovi le soluzioni dei 10 esercizi al PC di ogni modulo della dispensa,
    con il codice completo e l'\textbf{output reale} ottenuto eseguendolo.\\[0.4em]
    Aprilo \textbf{solo dopo} aver provato l'esercizio da solo: confrontare la tua
    soluzione con questa è il momento in cui si impara di più. Non esiste un'unica
    soluzione giusta: se la tua è diversa ma funziona e rispetta le regole, va bene.
  \end{tcolorbox}
  \vfill
  {\large a cura di \textbf{Giulio Casella}\par}
  \vspace{0.2cm}
  {\small Dispensa e Allenatore Prog2 -- materiale libero, da studenti per studenti\par}
\end{titlepage}

\tableofcontents

% Soluzioni degli esercizi del Modulo 1 (generato da separa_soluzioni.py)
\chapter{Soluzioni -- Modulo 1 -- Introduzione alla OOP \& Information Hiding}
\textbf{Esercizio 1.1}
\codice{codice/m01/es01.cpp}
\uscita{codice/m01/output/es01.txt}

\textbf{Esercizio 1.2}
\codice{codice/m01/es02.cpp}
\uscita{codice/m01/output/es02.txt}

\textbf{Esercizio 1.3}
\codice{codice/m01/es03.cpp}
\uscita{codice/m01/output/es03.txt}

\textbf{Esercizio 1.4} --- Memorizzare tre temperature vorrebbe dire tenerle
sempre allineate: basterebbe dimenticarsene una in un setter per avere un oggetto
incoerente. Con un solo dato le altre si calcolano e non possono essere sbagliate.
\codice{codice/m01/es04.cpp}
\uscita{codice/m01/output/es04.txt}

\textbf{Esercizio 1.5}
\codice{codice/m01/es05.cpp}
\uscita{codice/m01/output/es05.txt}

\textbf{Esercizio 1.6}
\codice{codice/m01/es06.cpp}
\uscita{codice/m01/output/es06.txt}

\textbf{Esercizio 1.7}
\codice{codice/m01/es07.cpp}
\uscita{codice/m01/output/es07.txt}

\textbf{Esercizio 1.8} --- Nota \cc|delta / 10.0|: con \cc|delta / 10| avresti una
divisione \textbf{intera} (45/10 = 4, non 4.5). L'output dipende dalla sequenza
di \cc|rand()| di MinGW: con un altro compilatore i numeri possono cambiare.
\codice{codice/m01/es08.cpp}
\uscita{codice/m01/output/es08.txt}

\textbf{Esercizio 1.9}
\codice{codice/m01/es09.cpp}
\uscita{codice/m01/output/es09.txt}

\textbf{Esercizio 1.10} --- \cc|mcd| è \cc|static|: non usa gli attributi di un
oggetto, è una funzione di servizio della classe. \cc|somma| e \cc|prodotto|
creano la nuova frazione con il costruttore, che la semplifica da solo:
l'invariante è garantita in un unico punto.
\codice{codice/m01/es10.cpp}
\uscita{codice/m01/output/es10.txt}

% Soluzioni degli esercizi del Modulo 2 (generato da separa_soluzioni.py)
\chapter{Soluzioni -- Modulo 2 -- Introduzione al C++ ed Elementi Base}
\textbf{Esercizio 2.1}
\codice{codice/m02/es01.cpp}
\ingresso{codice/m02/input/es01.txt}
\uscita{codice/m02/output/es01.txt}

\textbf{Esercizio 2.2} --- \cc|while (cin >> a >> op >> b)| si ferma da solo a
fine input o al primo dato illeggibile.
\codice{codice/m02/es02.cpp}
\ingresso{codice/m02/input/es02.txt}
\uscita{codice/m02/output/es02.txt}

\textbf{Esercizio 2.3}
\codice{codice/m02/es03.cpp}
\uscita{codice/m02/output/es03.txt}

\textbf{Esercizio 2.4} --- La riga A stampa 1 perché, nel momento in cui viene
eseguita, la \cc|x| locale non è ancora stata dichiarata: lo scope di una
variabile inizia dal punto della sua dichiarazione.
\codice{codice/m02/es04.cpp}
\uscita{codice/m02/output/es04.txt}

\textbf{Esercizio 2.5}
\codice{codice/m02/es05.cpp}
\ingresso{codice/m02/input/es05.txt}
\uscita{codice/m02/output/es05.txt}

\textbf{Esercizio 2.6}
\codice{codice/m02/es06.cpp}
\ingresso{codice/m02/input/es06.txt}
\uscita{codice/m02/output/es06.txt}

\textbf{Esercizio 2.7} --- Se leggessi direttamente \cc|double| dalla riga, al
primo \texttt{abc} lo stream andrebbe in errore e si fermerebbe tutto. Leggendo
\cc|string| la lettura non fallisce mai: la conversione avviene a parte, su uno
stream usa e getta.
\codice{codice/m02/es07.cpp}
\ingresso{codice/m02/input/es07.txt}
\uscita{codice/m02/output/es07.txt}

\textbf{Esercizio 2.8}
\codice{codice/m02/es08.cpp}
\uscita{codice/m02/output/es08.txt}

\textbf{Esercizio 2.9} --- tre file nella cartella \texttt{codice/m02/es09/}.
\codice{codice/m02/es09/geometria.h}
\codice{codice/m02/es09/geometria.cpp}
\codice{codice/m02/es09/main.cpp}
\uscita{codice/m02/output/es09.txt}

\textbf{Esercizio 2.10} --- La riga \texttt{questa non viene letta} si trova dopo
la riga vuota, che chiude lo scontrino.
\codice{codice/m02/es10.cpp}
\ingresso{codice/m02/input/es10.txt}
\uscita{codice/m02/output/es10.txt}

% Soluzioni degli esercizi del Modulo 3 (generato da separa_soluzioni.py)
\chapter{Soluzioni -- Modulo 3 -- Array, Puntatori, Reference, Stack e Heap}
\textbf{Esercizio 3.1} --- Gli indirizzi cambiano a ogni esecuzione: il sistema
operativo carica il programma ogni volta in una posizione diversa.
\codice{codice/m03/es01.cpp}
\uscita{codice/m03/output/es01.txt}

\textbf{Esercizio 3.2}
\codice{codice/m03/es02.cpp}
\uscita{codice/m03/output/es02.txt}

\textbf{Esercizio 3.3} --- \cc|pm - v| vale 5: la differenza tra due puntatori
nello stesso array è il numero di \emph{elementi} che li separa.
\codice{codice/m03/es03.cpp}
\uscita{codice/m03/output/es03.txt}

\textbf{Esercizio 3.4}
\codice{codice/m03/es04.cpp}
\uscita{codice/m03/output/es04.txt}

\textbf{Esercizio 3.5}
\codice{codice/m03/es05.cpp}
\uscita{codice/m03/output/es05.txt}

\textbf{Esercizio 3.6}
\codice{codice/m03/es06.cpp}
\ingresso{codice/m03/input/es06.txt}
\uscita{codice/m03/output/es06.txt}

\textbf{Esercizio 3.7} --- \cc|contiene(v, i, v[i])| cerca \cc|v[i]| solo tra gli
elementi \emph{precedenti}: se c'è già, è un duplicato.
\codice{codice/m03/es07.cpp}
\uscita{codice/m03/output/es07.txt}

\textbf{Esercizio 3.8}
\codice{codice/m03/es08.cpp}
\uscita{codice/m03/output/es08.txt}

\textbf{Esercizio 3.9}
\codice{codice/m03/es09.cpp}
\uscita{codice/m03/output/es09.txt}

\textbf{Esercizio 3.10} --- Per inserire si sposta partendo dal \textbf{fondo}
(\cc|i = n| scendendo), altrimenti ogni elemento sovrascriverebbe il successivo
prima di averlo spostato. Per rimuovere si parte invece dalla posizione tolta.
\codice{codice/m03/es10.cpp}
\uscita{codice/m03/output/es10.txt}

% Soluzioni degli esercizi del Modulo 4 (generato da separa_soluzioni.py)
\chapter{Soluzioni -- Modulo 4 -- Classi e Oggetti: Costruttori, Distruttore, Regola dei Tre}
\textbf{Esercizio 4.1}
\codice{codice/m04/es01.cpp}
\uscita{codice/m04/output/es01.txt}

\textbf{Esercizio 4.2} --- \texttt{GLOBALE} nasce prima di \texttt{inizio main} e
muore dopo \texttt{fine main}. \texttt{HEAP} muore quando lo decide il \cc|delete|,
non alla fine del blocco.
\codice{codice/m04/es02.cpp}
\uscita{codice/m04/output/es02.txt}

\textbf{Esercizio 4.3} --- Dentro \cc|perValore| i vivi sono 3: il parametro è una
copia, creata con il costruttore di copia e distrutta all'uscita.
\codice{codice/m04/es03.cpp}
\uscita{codice/m04/output/es03.txt}

\textbf{Esercizio 4.4}
\codice{codice/m04/es04.cpp}
\uscita{codice/m04/output/es04.txt}

\textbf{Esercizio 4.5} --- \cc|Punto2D(int x = 0, int y = 0) : x(x), y(y)|
funziona: nella lista, fuori dalle parentesi c'è l'attributo, dentro il parametro.
\codice{codice/m04/es05.cpp}
\uscita{codice/m04/output/es05.txt}

\textbf{Esercizio 4.6}
\codice{codice/m04/es06.cpp}
\uscita{codice/m04/output/es06.txt}

\textbf{Esercizio 4.7}
\codice{codice/m04/es07.cpp}
\uscita{codice/m04/output/es07.txt}

\textbf{Esercizio 4.8} --- In \cc|concatena| la stringa temporanea \cc|tmp| va
liberata, perché il costruttore di \cc|MiaStringa| ne fa una copia sua.
\codice{codice/m04/es08.cpp}
\uscita{codice/m04/output/es08.txt}

\textbf{Esercizio 4.9} --- I metodi privati \cc|allocaMemoria| e \cc|liberaMemoria|
evitano di ripetere lo stesso codice in costruttori, \cc|operator=| e distruttore.
\codice{codice/m04/es09.cpp}
\uscita{codice/m04/output/es09.txt}

\textbf{Esercizio 4.10} --- \cc|explicit| davanti al costruttore impedisce
conversioni automatiche indesiderate, come \cc|Pila p = 5;|.
\codice{codice/m04/es10.cpp}
\uscita{codice/m04/output/es10.txt}

% Soluzioni degli esercizi del Modulo 5 (generato da separa_soluzioni.py)
\chapter{Soluzioni -- Modulo 5 -- Template di Funzione e di Classe}
\textbf{Esercizio 5.1}
\codice{codice/m05/es01.cpp}
\uscita{codice/m05/output/es01.txt}

\textbf{Esercizio 5.2}
\codice{codice/m05/es02.cpp}
\uscita{codice/m05/output/es02.txt}

\textbf{Esercizio 5.3} --- \cc|T()| è il ``valore zero'' del tipo: 0 per i numeri,
stringa vuota per \cc|string|. Per le stringhe \cc|+=| concatena.
\codice{codice/m05/es03.cpp}
\uscita{codice/m05/output/es03.txt}

\textbf{Esercizio 5.4} --- Per contare \q{rosso} serve \cc|string("rosso")|:
con il letterale \cc|"rosso"| il terzo argomento sarebbe un \cc|const char*| e
\cc|T| risulterebbe dedotto in due modi diversi.
\codice{codice/m05/es04.cpp}
\uscita{codice/m05/output/es04.txt}

\textbf{Esercizio 5.5}
\codice{codice/m05/es05.cpp}
\ingresso{codice/m05/input/es05.txt}
\uscita{codice/m05/output/es05.txt}

\textbf{Esercizio 5.6}
\codice{codice/m05/es06.cpp}
\uscita{codice/m05/output/es06.txt}

\textbf{Esercizio 5.7} --- \cc|invertita()| restituisce un tipo diverso da quello
dell'oggetto (\cc|Coppia<U, T>|): dentro un template si possono usare altre
istanziazioni dello stesso template.
\codice{codice/m05/es07.cpp}
\uscita{codice/m05/output/es07.txt}

\textbf{Esercizio 5.8} --- Attenzione a
\cc|cout << lettere.estrarre() << lettere.estrarre()|: dal C++17 le espressioni
concatenate con \cc|<<| sono valutate da sinistra a destra, quindi esce prima
\texttt{b} e poi \texttt{a}. Con compilatori più vecchi l'ordine non era garantito:
meglio usare due istruzioni separate.
\codice{codice/m05/es08.cpp}
\uscita{codice/m05/output/es08.txt}

\textbf{Esercizio 5.9}
\codice{codice/m05/es09.cpp}
\uscita{codice/m05/output/es09.txt}

\textbf{Esercizio 5.10}
\codice{codice/m05/es10.cpp}
\uscita{codice/m05/output/es10.txt}

% Soluzioni degli esercizi del Modulo 6 (generato da separa_soluzioni.py)
\chapter{Soluzioni -- Modulo 6 -- Complessità Computazionale}
\textbf{Esercizio 6.1} --- Il rapporto resta fermo a 0.5: i passi crescono come $n$,
quindi $O(n)$.
\codice{codice/m06/es01.cpp}
\uscita{codice/m06/output/es01.txt}

\textbf{Esercizio 6.2} --- Il triangolo fa circa metà dei passi, ma il rapporto con
$n^2$ si stabilizza: entrambi $O(n^2)$.
\codice{codice/m06/es02.cpp}
\uscita{codice/m06/output/es02.txt}

\textbf{Esercizio 6.3} --- $n \times 10$ significa solo $+3$ o $+4$ passi:
logaritmico.
\codice{codice/m06/es03.cpp}
\uscita{codice/m06/output/es03.txt}

\textbf{Esercizio 6.4}
\codice{codice/m06/es04.cpp}
\uscita{codice/m06/output/es04.txt}

\textbf{Esercizio 6.5} --- Caso medio $(n+1)/2 = 500.5$: sempre $O(n)$.
\codice{codice/m06/es05.cpp}
\uscita{codice/m06/output/es05.txt}

\textbf{Esercizio 6.6} --- $O(n \log^2 n)$, cioè $n \cdot \log n \cdot \log n$.
\codice{codice/m06/es06.cpp}
\uscita{codice/m06/output/es06.txt}

\textbf{Esercizio 6.7}
\codice{codice/m06/es07.cpp}
\uscita{codice/m06/output/es07.txt}

\textbf{Esercizio 6.8} --- Se $n = a \cdot b$ con $a \le b$, allora $a \cdot a \le n$:
il divisore più piccolo, se esiste, è al più la radice di $n$. Da $O(n)$ a
$O(\sqrt{n})$.
\codice{codice/m06/es08.cpp}
\uscita{codice/m06/output/es08.txt}

\textbf{Esercizio 6.9} --- Un milione di confronti contro 19.
\codice{codice/m06/es09.cpp}
\uscita{codice/m06/output/es09.txt}

\textbf{Esercizio 6.10} --- A ogni $+5$ le chiamate si moltiplicano per circa 11:
crescita esponenziale. Il ciclo invece cresce di 5 giri.
\codice{codice/m06/es10.cpp}
\uscita{codice/m06/output/es10.txt}

% Soluzioni degli esercizi del Modulo 7 (generato da separa_soluzioni.py)
\chapter{Soluzioni -- Modulo 7 -- Algoritmi di Ordinamento e Ricerca}
\textbf{Esercizio 7.1}
\codice{codice/m07/es01.cpp}
\uscita{codice/m07/output/es01.txt}

\textbf{Esercizio 7.2}
\codice{codice/m07/es02.cpp}
\uscita{codice/m07/output/es02.txt}

\textbf{Esercizio 7.3} --- Nel codice ASCII le maiuscole (65--90) vengono prima delle
minuscole (97--122), quindi \q{bianchi} va dopo \q{Verdi}.
\codice{codice/m07/es03.cpp}
\uscita{codice/m07/output/es03.txt}

\textbf{Esercizio 7.4} --- Il risultato si salva in una variabile \emph{prima} di
stamparlo: se la chiamata fosse dentro lo stesso \cc|cout|, la scritta
\texttt{-> indice} comparirebbe prima dei passi.
\codice{codice/m07/es04.cpp}
\uscita{codice/m07/output/es04.txt}

\textbf{Esercizio 7.5}
\codice{codice/m07/es05.cpp}
\uscita{codice/m07/output/es05.txt}

\textbf{Esercizio 7.6} --- Con \cc|<|, a parità di voto verrebbe preso prima
l'elemento della metà destra, e l'ordine alfabetico si romperebbe.
\codice{codice/m07/es06.cpp}
\uscita{codice/m07/output/es06.txt}

\textbf{Esercizio 7.7}
\codice{codice/m07/es07.cpp}
\uscita{codice/m07/output/es07.txt}

\textbf{Esercizio 7.8} --- $O(n + k)$, con $k$ = numero di valori possibili (13).
\codice{codice/m07/es08.cpp}
\uscita{codice/m07/output/es08.txt}

\textbf{Esercizio 7.9}
\codice{codice/m07/es09.cpp}
\uscita{codice/m07/output/es09.txt}

\textbf{Esercizio 7.10}
\codice{codice/m07/es10.cpp}
\uscita{codice/m07/output/es10.txt}

% Soluzioni degli esercizi del Modulo 8
\chapter{Soluzioni -- Modulo 8 -- Ricorsione}

\textbf{Esercizio 8.1} --- \cc|suEGiu(3)| stampa all'andata 3, 2, 1 e al ritorno
1, 2, 3: ogni numero compare due volte, ``a specchio''.
\codice{codice/m08/es01.cpp}
\uscita{codice/m08/output/es01.txt}

\textbf{Esercizio 8.2} --- Con caso base \cc|b == 1| la chiamata \cc|prod(5, 0)| non
termina: \cc|b| scende a $-1, -2, \ldots$ senza mai valere 1.
\codice{codice/m08/es02.cpp}
\uscita{codice/m08/output/es02.txt}

\textbf{Esercizio 8.3}
\codice{codice/m08/es03.cpp}
\uscita{codice/m08/output/es03.txt}

\textbf{Esercizio 8.4} --- Le righe \texttt{->} compaiono in ordine inverso rispetto
alle chiamate: è il ritorno.
\codice{codice/m08/es04.cpp}
\uscita{codice/m08/output/es04.txt}

\textbf{Esercizio 8.5}
\codice{codice/m08/es05.cpp}
\uscita{codice/m08/output/es05.txt}

\textbf{Esercizio 8.6} --- In \cc|contrario|, \cc|contrario(s.substr(1)) + s[0]|
mette il primo carattere \emph{in fondo} al contrario del resto.
\codice{codice/m08/es06.cpp}
\uscita{codice/m08/output/es06.txt}

\textbf{Esercizio 8.7} --- È di coda: la chiamata è l'ultima cosa. Nel ciclo, i
parametri \cc|a| e \cc|b| diventano variabili aggiornate a ogni giro.
\codice{codice/m08/es07.cpp}
\uscita{codice/m08/output/es07.txt}

\textbf{Esercizio 8.8} --- 27 fa ben 111 passi, salendo fino a 9232 prima di
scendere. Che ogni numero arrivi a 1 non è mai stato dimostrato.
\codice{codice/m08/es08.cpp}
\uscita{codice/m08/output/es08.txt}

\textbf{Esercizio 8.9} --- Le mosse sono $2^n - 1$: per 20 dischi più di un milione.
\codice{codice/m08/es09.cpp}
\uscita{codice/m08/output/es09.txt}

\textbf{Esercizio 8.10} --- Il secondo \cc|swap|, dopo la chiamata, rimette la stringa
com'era: senza, i tentativi successivi partirebbero da una stringa già modificata.
\codice{codice/m08/es10.cpp}
\uscita{codice/m08/output/es10.txt}

% Soluzioni degli esercizi del Modulo 9
\chapter{Soluzioni -- Modulo 9 -- Ereditarietà, Polimorfismo, Classi Astratte e RTTI}

\textbf{Esercizio 9.1} --- Costruttori dalla base alla derivata, distruttori al
contrario, alla chiusura del blocco.
\codice{codice/m09/es01.cpp}
\uscita{codice/m09/output/es01.txt}

\textbf{Esercizio 9.2} --- Le righe commentate sono quelle che non compilano.
\codice{codice/m09/es02.cpp}
\uscita{codice/m09/output/es02.txt}

\textbf{Esercizio 9.3} --- Tramite \cc|Veicolo*| il metodo non virtuale esegue la
versione di \cc|Veicolo|, quello virtuale la versione di \cc|Moto|. Sulla copia
\cc|Veicolo copia = m;| anche il virtuale esegue la versione di \cc|Veicolo|: la copia
\emph{è} un \cc|Veicolo| (slicing).
\codice{codice/m09/es03.cpp}
\uscita{codice/m09/output/es03.txt}

\textbf{Esercizio 9.4}
\codice{codice/m09/es04.cpp}
\uscita{codice/m09/output/es04.txt}

\textbf{Esercizio 9.5} --- \cc|Color(c)| converte l'intero letto in un valore
dell'\cc|enum|. L'array di puntatori è dinamico perché la sua dimensione si conosce solo
leggendo l'input: si libera prima ogni figura, poi l'array.
\codice{codice/m09/es05.cpp}
\ingresso{codice/m09/input/es05.txt}
\uscita{codice/m09/output/es05.txt}

\textbf{Esercizio 9.6} --- Senza \cc|virtual| il distruttore della derivata non parte e
il suo buffer non viene liberato.
\codice{codice/m09/es06.cpp}
\uscita{codice/m09/output/es06.txt}

\textbf{Esercizio 9.7} --- \cc|typeid(...).name()| restituisce \texttt{4Auto} con g++:
un nome interno del compilatore, utile solo per il debug.
\codice{codice/m09/es07.cpp}
\uscita{codice/m09/output/es07.txt}

\textbf{Esercizio 9.8} --- \cc|to_string| converte il numero in testo; \cc|c_str()|
restituisce la stringa nel formato \cc|const char*| che \cc|what()| richiede. Il messaggio
è salvato in un attributo, così resta valido anche dopo la fine del costruttore.
\codice{codice/m09/es08.cpp}
\uscita{codice/m09/output/es08.txt}

\textbf{Esercizio 9.9} --- Si ordinano i \emph{puntatori}: gli oggetti restano dove
sono nello heap, cambia solo l'ordine delle frecce.
\codice{codice/m09/es09.cpp}
\uscita{codice/m09/output/es09.txt}

\textbf{Esercizio 9.10} --- Senza \cc|virtual| il costruttore di \cc|A| parte due volte:
in \cc|D| ci sono due \cc|A|. Con l'ereditarietà virtuale una volta sola.
\codice{codice/m09/es10.cpp}
\uscita{codice/m09/output/es10.txt}

% Soluzioni degli esercizi del Modulo 10
\chapter{Soluzioni -- Modulo 10 -- Sovraccaricamento degli Operatori}

\textbf{Esercizio 10.1}
\codice{codice/m10/es01.cpp}
\uscita{codice/m10/output/es01.txt}

\textbf{Esercizio 10.2} --- Memorizzare un solo numero (i secondi totali) rende la
somma banale; la conversione in ore, minuti e secondi si fa solo nella stampa.
\codice{codice/m10/es02.cpp}
\uscita{codice/m10/output/es02.txt}

\textbf{Esercizio 10.3} --- Gli operatori esterni non hanno bisogno di \cc|friend| se la
classe offre getter pubblici sufficienti.
\codice{codice/m10/es03.cpp}
\uscita{codice/m10/output/es03.txt}

\textbf{Esercizio 10.4} --- \cc|++(++c)| funziona perché il prefisso restituisce
\cc|Contatore&|: il secondo \cc|++| agisce sullo stesso oggetto.
\codice{codice/m10/es04.cpp}
\uscita{codice/m10/output/es04.txt}

\textbf{Esercizio 10.5} --- Il costruttore porta il segno al numeratore (\cc|5/-6|
diventa \cc|-5/6|), così nel confronto incrociato i denominatori sono positivi e la
disuguaglianza non si capovolge.
\codice{codice/m10/es05.cpp}
\uscita{codice/m10/output/es05.txt}

\textbf{Esercizio 10.6}
\codice{codice/m10/es06.cpp}
\uscita{codice/m10/output/es06.txt}

\textbf{Esercizio 10.7} --- Dopo un fallimento, \cc|clear()| e \cc|ignore()| (Modulo 2)
saltano la riga sbagliata. \texttt{Bruno 33} viene letto correttamente come testo e
numero, ma scartato per il valore fuori intervallo.
\codice{codice/m10/es07.cpp}
\ingresso{codice/m10/input/es07.txt}
\uscita{codice/m10/output/es07.txt}

\textbf{Esercizio 10.8}
\codice{codice/m10/es08.cpp}
\uscita{codice/m10/output/es08.txt}

\textbf{Esercizio 10.9} --- \cc|+=| fa tutto il lavoro di inserimento (controllo dei
duplicati, crescita): unione, intersezione e differenza lo riusano.
\codice{codice/m10/es09.cpp}
\uscita{codice/m10/output/es09.txt}

\textbf{Esercizio 10.10} --- \cc|operator-| conta i giorni applicando \cc|++| finché le
date coincidono: semplice e sicuro, anche se lento per intervalli enormi.
\codice{codice/m10/es10.cpp}
\uscita{codice/m10/output/es10.txt}

% Soluzioni degli esercizi del Modulo 11
\chapter{Soluzioni -- Modulo 11 -- Liste Concatenate}

\textbf{Esercizio 11.1} --- \cc|new Nodo{5, primo}| crea il nodo inizializzando i campi
nell'ordine in cui sono dichiarati nella \cc|struct|.
\codice{codice/m11/es01.cpp}
\uscita{codice/m11/output/es01.txt}

\textbf{Esercizio 11.2}
\codice{codice/m11/es02.cpp}
\uscita{codice/m11/output/es02.txt}

\textbf{Esercizio 11.3}
\codice{codice/m11/es03.cpp}
\uscita{codice/m11/output/es03.txt}

\textbf{Esercizio 11.4} --- La testa va passata per reference perché
\cc|rimuoviTesta| (e \cc|rimuoviCoda| con un solo nodo) deve cambiare il puntatore del
chiamante, non una sua copia.
\codice{codice/m11/es04.cpp}
\uscita{codice/m11/output/es04.txt}

\textbf{Esercizio 11.5} --- In \cc|rimuovi| il ciclo si ferma al primo valore non minore:
in una lista ordinata, se lì non c'è il valore cercato, non c'è nemmeno più avanti.
\codice{codice/m11/es05.cpp}
\uscita{codice/m11/output/es05.txt}

\textbf{Esercizio 11.6} --- Nella versione ricorsiva, \cc|testa->succ| è (prima della
chiamata) il secondo nodo, che dopo l'inversione del resto è diventato l'ultimo: gli si
aggancia la testa.
\codice{codice/m11/es06.cpp}
\uscita{codice/m11/output/es06.txt}

\textbf{Esercizio 11.7} --- Nell'inserimento prima di \cc|pos| si aggiornano quattro
frecce: \cc|nuovo->succ|, \cc|nuovo->prec|, il \cc|succ| del vecchio precedente (o la
testa) e \cc|pos->prec|.
\codice{codice/m11/es07.cpp}
\uscita{codice/m11/output/es07.txt}

\textbf{Esercizio 11.8} --- Il nodo \cc|finto| è una variabile locale: non va
cancellato, e la lista risultato parte da \cc|finto.succ|.
\codice{codice/m11/es08.cpp}
\uscita{codice/m11/output/es08.txt}

\textbf{Esercizio 11.9}
\codice{codice/m11/es09.cpp}
\uscita{codice/m11/output/es09.txt}

\textbf{Esercizio 11.10} --- Nessun \cc|if| per la lista vuota o per la testa: è il
vantaggio della sentinella. Il costruttore usa \cc|T()| per il valore (inutile) della
sentinella.
\codice{codice/m11/es10.cpp}
\uscita{codice/m11/output/es10.txt}

% Soluzioni degli esercizi del Modulo 12
\chapter{Soluzioni -- Modulo 12 -- Pile e Code}

\textbf{Esercizio 12.1} --- La capacità è un parametro template: \cc|PilaArray<int, 3>|
usa un array statico di 3 elementi.
\codice{codice/m12/es01.cpp}
\uscita{codice/m12/output/es01.txt}

\textbf{Esercizio 12.2}
\codice{codice/m12/es02.cpp}
\uscita{codice/m12/output/es02.txt}

\textbf{Esercizio 12.3} --- \cc|(pagine + 19) / 20| è la divisione intera arrotondata
per eccesso. Il biglietto da una pagina aspetta 6 minuti: in una coda FIFO conta solo
l'ordine di arrivo (l'Esercizio~12.10 mostra l'alternativa con priorità).
\codice{codice/m12/es03.cpp}
\uscita{codice/m12/output/es03.txt}

\textbf{Esercizio 12.4}
\codice{codice/m12/es04.cpp}
\uscita{codice/m12/output/es04.txt}

\textbf{Esercizio 12.5}
\codice{codice/m12/es05.cpp}
\uscita{codice/m12/output/es05.txt}

\textbf{Esercizio 12.6} --- Ogni elemento viene spostato al più una volta da
\cc|entrata| a \cc|uscita|: in media ogni operazione costa $O(1)$.
\codice{codice/m12/es06.cpp}
\uscita{codice/m12/output/es06.txt}

\textbf{Esercizio 12.7}
\codice{codice/m12/es07.cpp}
\uscita{codice/m12/output/es07.txt}

\textbf{Esercizio 12.8} --- Le pile memorizzano versioni intere del testo: semplice ma
costoso per testi lunghi. Un editor vero memorizzerebbe le singole operazioni.
\codice{codice/m12/es08.cpp}
\uscita{codice/m12/output/es08.txt}

\textbf{Esercizio 12.9} --- \texttt{8 - 2 - 1} diventa \texttt{8 2 - 1 -}: a parità di
priorità esce prima l'operatore già in pila, così la sottrazione resta associativa a
sinistra ($(8 - 2) - 1 = 5$).
\codice{codice/m12/es09.cpp}
\uscita{codice/m12/output/es09.txt}

\textbf{Esercizio 12.10} --- L'inserimento si ferma \emph{dopo} tutti i pazienti con
codice minore o uguale: per questo a parità di codice resta l'ordine di arrivo.
\codice{codice/m12/es10.cpp}
\uscita{codice/m12/output/es10.txt}

% Soluzioni degli esercizi del Modulo 13
\chapter{Soluzioni -- Modulo 13 -- Alberi Binari e Alberi Binari di Ricerca (BST)}

\textbf{Esercizio 13.1}
\codice{codice/m13/es01.cpp}
\uscita{codice/m13/output/es01.txt}

\textbf{Esercizio 13.2} --- Lo schema ``caso base sull'albero vuoto, poi io + sinistra +
destra'' risolve quasi tutti i problemi di conteggio.
\codice{codice/m13/es02.cpp}
\uscita{codice/m13/output/es02.txt}

\textbf{Esercizio 13.3}
\codice{codice/m13/es03.cpp}
\uscita{codice/m13/output/es03.txt}

\textbf{Esercizio 13.4}
\codice{codice/m13/es04.cpp}
\uscita{codice/m13/output/es04.txt}

\textbf{Esercizio 13.5} --- \cc|Nodo** p| è un puntatore al \emph{posto} in cui andrà
il nuovo nodo (la radice o un campo \cc|sx|/\cc|dx|): così l'inserimento iterativo non ha
bisogno di casi speciali. Il \cc|distruggi| usa un ciclo sul ramo destro per non
annidare 1000 chiamate.
\codice{codice/m13/es05.cpp}
\uscita{codice/m13/output/es05.txt}

\textbf{Esercizio 13.6}
\codice{codice/m13/es06.cpp}
\uscita{codice/m13/output/es06.txt}

\textbf{Esercizio 13.7}
\codice{codice/m13/es07.cpp}
\uscita{codice/m13/output/es07.txt}

\textbf{Esercizio 13.8} --- \cc|LONG_MIN| e \cc|LONG_MAX| (da \cc|<climits>|) fanno da
``nessun limite'' iniziale.
\codice{codice/m13/es08.cpp}
\uscita{codice/m13/output/es08.txt}

\textbf{Esercizio 13.9}
\codice{codice/m13/es09.cpp}
\uscita{codice/m13/output/es09.txt}

\textbf{Esercizio 13.10} --- \cc|to_string| converte il valore in testo, così si può usare
\q{-} quando il successore o il predecessore non esistono.
\codice{codice/m13/es10.cpp}
\uscita{codice/m13/output/es10.txt}

% Soluzioni degli esercizi del Modulo 14
\chapter{Soluzioni -- Modulo 14 -- Grafi}

\textbf{Esercizio 14.1} --- Il numero di archi è la somma dei gradi divisa per 2, perché
ogni arco viene contato da entrambi i suoi estremi.
\codice{codice/m14/es01.cpp}
\uscita{codice/m14/output/es01.txt}

\textbf{Esercizio 14.2} --- Con l'inserimento in testa i vicini compaiono in ordine
inverso rispetto a quello di inserimento.
\codice{codice/m14/es02.cpp}
\uscita{codice/m14/output/es02.txt}

\textbf{Esercizio 14.3} --- Per avere le liste in ordine crescente inserendo in testa,
basta scorrere ogni riga della matrice da destra a sinistra: l'ultimo inserito (il più
piccolo) finisce in testa.
\codice{codice/m14/es03.cpp}
\uscita{codice/m14/output/es03.txt}

\textbf{Esercizio 14.4} --- Il grado uscente di $u$ è la lunghezza della sua lista; il
grado entrante è il numero di volte in cui $u$ compare nelle liste degli \emph{altri}.
Il trasposto si costruisce con un solo passaggio: $O(|V| + |E|)$.
\codice{codice/m14/es04.cpp}
\uscita{codice/m14/output/es04.txt}

\textbf{Esercizio 14.5} --- Qui \cc|dist[v] == -1| fa anche da \q{bianco}: non serve un
array di colori separato.
\codice{codice/m14/es05.cpp}
\uscita{codice/m14/output/es05.txt}

\textbf{Esercizio 14.6} --- Nella DFS iterativa un vertice può finire in pila più volte
(lo mettono due vicini diversi prima che venga estratto): per questo si controlla
\cc|vis[u]| quando lo si \emph{estrae}. L'ordine di visita (0 2 3 4 1) è diverso da quello
ricorsivo, perché la pila restituisce per ultimo il primo vicino inserito.
\codice{codice/m14/es06.cpp}
\uscita{codice/m14/output/es06.txt}

\textbf{Esercizio 14.7}
\codice{codice/m14/es07.cpp}
\uscita{codice/m14/output/es07.txt}

\textbf{Esercizio 14.8} --- Con due soli colori (visitato sì/no), nel rombo il vertice 3
viene raggiunto una seconda volta da 2 dopo essere stato visitato da 1: sembrerebbe un
ciclo, ma 3 è già \textbf{finito} (nero), non è un antenato aperto. Solo un arco verso un
vertice grigio chiude davvero un ciclo.
\codice{codice/m14/es08.cpp}
\uscita{codice/m14/output/es08.txt}

\textbf{Esercizio 14.9} --- Il grafo è \emph{implicito}: non c'è né matrice né lista, i
quattro vicini di una cella si calcolano con gli spostamenti \cc|dr[]| e \cc|dc[]|. La
BFS garantisce che il percorso trovato sia il più corto (19 passi).
\codice{codice/m14/es09.cpp}
\uscita{codice/m14/output/es09.txt}

\textbf{Esercizio 14.10} --- Un grafo è bipartito se e solo se non contiene cicli di
lunghezza dispari: l'esagono (ciclo di 6) lo è, il triangolo (ciclo di 3) no. La coda è
un oggetto locale: uscendo con \cc|return false| a metà, il suo distruttore libera i nodi
rimasti.
\codice{codice/m14/es10.cpp}
\uscita{codice/m14/output/es10.txt}


\appendix
% Soluzioni della simulazione d'esame (Appendice A.2 della dispensa)
\chapter{Soluzioni -- Simulazione d'esame}

\section*{Prova 1 -- Quiz a crocette}

Conta le risposte giuste: con 16 o più sei in buona forma; sotto le 12, ripassa i moduli
delle domande sbagliate (indicati tra parentesi).

\begin{soluzione}[: risposte commentate]
\begin{enumerate}[label=\textbf{\arabic*.}]
  \item \textbf{B} (M1).
  \item \textbf{A} (M2). La \texttt{x} interna è un'altra variabile che \q{nasconde} quella
        esterna solo dentro il blocco; fuori, \texttt{x} vale ancora 5.
  \item \textbf{C} (M3). Aritmetica dei puntatori: \texttt{p + 2} punta ad \texttt{a[2]},
        l'asterisco ne legge il valore.
  \item \textbf{B} (M3). La A è un dangling pointer.
  \item \textbf{C} (M4). Regola dei Tre.
  \item \textbf{B} (M4). Gli oggetti locali vengono distrutti in ordine inverso di creazione.
  \item \textbf{B} (M5).
  \item \textbf{D} (M6). Si tiene il termine che cresce di più: $n^2$ cresce più di
        $n \log n$.
  \item \textbf{B} (M6). Il ciclo interno fa $0 + 1 + \ldots + (n-1)$ giri, circa
        $n^2/2$.
  \item \textbf{A} (M7).
  \item \textbf{C} (M7). Il quick sort è $O(n^2)$ nel caso peggiore.
  \item \textbf{C} (M8).
  \item \textbf{B} (M9). Senza \texttt{virtual} il binding è statico: decide il tipo del
        puntatore, \texttt{Base*}.
  \item \textbf{A} (M9). Può però avere costruttori, attributi, ed essere tipo di puntatori.
  \item \textbf{B} (M10). In \texttt{cout << x} l'operando sinistro è \texttt{cout}: un
        metodo dovrebbe appartenere a \texttt{ostream}.
  \item \textbf{A} (M11).
  \item \textbf{B} (M12). Dopo i tre push la pila è 1 2 3 (cima 3): pop dà 3, pop dà 2;
        push(4) mette 4 in cima, pop dà 4.
  \item \textbf{D} (M13).
  \item \textbf{B} (M14).
  \item \textbf{C} (M10). Così si possono concatenare gli assegnamenti
        (\texttt{a = b = c}) senza copie inutili.
\end{enumerate}
\end{soluzione}

\section*{Prova 2 -- Laboratorio: una soluzione completa}

La soluzione segue i punti della traccia con commenti \texttt{// ===== STEP n =====}. Non
è l'unica possibile: conta che ogni punto funzioni e sia mostrato nell'output.
\codice{codice/esame/lab.cpp}
\uscita{codice/esame/output/lab.txt}

\begin{consiglio}[: cosa notare nella soluzione]
\begin{itemize}
  \item \texttt{typeid(...).name()} con MinGW stampa il nome \q{decorato} (\texttt{4Auto}:
        la cifra è la lunghezza del nome). Su altri compilatori può essere diverso: è
        normale.
  \item Coda e BST contengono \texttt{Veicolo*}: i loro distruttori liberano i
        \textbf{nodi}, non i veicoli. I veicoli li ha creati il \texttt{main}, ed è il
        \texttt{main} a liberarli alla fine. Se anche le strutture li cancellassero, ogni
        veicolo verrebbe cancellato due volte.
  \item Nel BST i tre tagliandi da 160 euro hanno la stessa chiave: vanno a destra (regola
        dei duplicati del Modulo 13) e la visita in-order li stampa nell'ordine di
        inserimento.
  \item L'eccezione viene lanciata dentro \texttt{dequeue} e catturata nel \texttt{main}
        \textbf{per riferimento costante}; l'incasso calcolato fino a quel momento resta
        valido.
  \item \texttt{typeid} confronta il tipo \emph{esatto}, \texttt{dynamic\_cast} permette
        anche di usare i metodi propri della derivata (\texttt{getCilindrata}).
  \item Le scelte di progetto discutibili (l'eccezione \texttt{CodaVuota} riusata in
        \texttt{massimo} del BST, le statistiche calcolate nel \texttt{main}) vanno sapute
        giustificare all'orale: \q{per semplicità} è una risposta accettabile, se sai dire
        come faresti meglio (un'eccezione \texttt{AlberoVuoto}, un metodo dedicato).
\end{itemize}
\end{consiglio}

\section*{Prova 3 -- Orale}

Le risposte modello sono nei moduli indicati nella tabella della simulazione. Valuta te
stesso così: hai risposto per almeno un minuto senza fermarti? Hai usato i termini
giusti (glossario)? Hai fatto un esempio o un disegno? Se una delle tre manca, ripeti la
domanda il giorno dopo.

